\chapter{测度、归一化与基础工具}
\label{app:measure}
\section{变长轨迹如何诱导条件目标}
设提示 $x\sim\mathcal D$，锚点 $q$ 生成轨迹，回复长度 $H$ 包含 EOS，不包含提示与 padding。预算终止时保留已生成动作。假设 $0<\E H<\infty$，定义
\begin{equation}
d_q(B)=\frac{\E_q\sum_{t\ge1}\mathbf1\{t\le H,\ s_t\in B\}}{\E_qH}.
\end{equation}
\begin{proposition}[访问测度的条件分解]
在上述停止规则及 $0<\E_qH<\infty$ 下，总体前缀--动作测度为 $d_q(s)q(a\mid s)$。对满足 $\E_q\sum_{t\le H}|h(s_t,a_t)|<\infty$ 的函数，
\begin{equation}
\E_q\sum_{t\le H}h(s_t,a_t)=\E_qH\;\E_{d_q,q}h.
\end{equation}
\end{proposition}
\begin{proof}
这里 $s_t$ 是抽取第 $t$ 个动作之前的历史。存活事件 $\{t\le H\}$ 关于这段历史可测。对于绝对可积的 $h$，塔式期望给出
\begin{align}
\E_q\sum_{t\ge1}\mathbf1\{t\le H\}h(s_t,a_t)
&=\E_q\sum_{t\ge1}\mathbf1\{t\le H\}
\sum_aq(a\mid s_t)h(s_t,a)\\
&=\E_qH\ \E_{s\sim d_q,a\sim q(\cdot\mid s)}h(s,a).
\end{align}
第一步交换求和与期望由非负性或绝对可积性保证。因此总体前缀--动作测度是 $d_qq$。
\end{proof}

\begin{proposition}[全局 token 比率估计的一致性]
独立同分布轨迹具有有限正平均长度及可积的轨迹总损失时，全局 token 均值
$\widehat I_B=(\sum_{i,t\le H_i}h_{it})/(\sum_iH_i)$ 几乎处处收敛到 $\E_{d_q,q}h$。该结论不要求轨迹内 token 独立，也不保证有限样本无偏。
\end{proposition}
\begin{proof}
对 $B$ 条独立轨迹，全局 token 均值为
\begin{equation}
\widehat I_B=\frac{\sum_{i=1}^B\sum_{t=1}^{H_i}h(s_{it},a_{it})}{\sum_{i=1}^BH_i}.
\end{equation}
分子分母除以 $B$，分别由大数定律趋于对应期望；分母极限严格为正，所以比值一致收敛到 $\E_{d_q,q}h$。有限 $B$ 时随机分母一般产生比率估计偏差。同一轨迹内的 token 相关，置信区间不能把全部 token 当成独立样本。
\end{proof}

\begin{proposition}[逐序列长度权重改变动作分布]
存在有限长度生成过程，使按每条轨迹的 $1/H$ 加权后，同一前缀下的动作分布不同于生成策略。因此逐序列均值不能一般替代 $d_qq$ 下的 token 目标。
\end{proposition}
\begin{proof}
逐序列先除以 $H_i$ 会改变测度。取初始两动作各以 $1/2$ 发生，动作一立即终止，动作二随后再生成一个确定 EOS。按 $1/H$ 加权后，初始两动作质量为 $1/2,1/4$，条件概率成为 $2/3,1/3$。这个权重依赖当前与未来动作，不能移出条件期望。
\end{proof}

\begin{proposition}[微批次梯度累加不变性]
固定有效 token 总数 $\mathsf T>0$，微批次构成有效位置的不交划分。若每批使用总分母 $\mathsf T$，并在全部梯度累加后进行裁剪和更新，则所得梯度等于全 batch 目标的梯度。
\end{proposition}
\begin{proof}
固定 batch 后，总有效 token 数 $\mathsf T$ 不依赖内层参数。若微批次互不相交，使用 $L_m=\mathsf T^{-1}\sum_{i\in\mathcal B_m}\ell_i$，则
\begin{equation}
\sum_m\nabla L_m=\mathsf T^{-1}\sum_i\nabla\ell_i.
\end{equation}
每个微批次各自取均值时，需再乘 $|\mathcal B_m|/\mathsf T$。裁剪和优化器更新放在梯度累加以后，才保留这条线性恒等式。
\end{proof}

\section{积分微分与支持条件}
默认有限词表上的概率严格为正。若允许零概率，$\KL(P\Vert Q)$ 有限至少要求 $P\ll Q$；重要性抽样则需要提案覆盖非零被积贡献。两项支持要求由不同对象决定。

\begin{lemma}[score 恒等式与积分微分]
有限词表上，若 $p_\theta>0$ 归一化且 $p_\theta,h_\theta$ 可微，则
\begin{equation}
\E_p\psi=0,\qquad
\nabla_\theta\E_ph_\theta=\E_p(\nabla_\theta h_\theta+h_\theta\psi).
\end{equation}
若积分权重 $q$ 冻结，则 $\nabla_\theta\E_qh_\theta=\E_q\nabla_\theta h_\theta$。
\end{lemma}
\begin{proof}
令 $p_\theta$ 可微，$\psi_a=\nabla_\theta\log p_\theta(a)$。归一化给出
\begin{equation}
\E_p\psi=\sum_a\nabla p_\theta(a)=0.
\end{equation}
对参数依赖函数 $h_\theta$，
\begin{equation}
\nabla_\theta\E_{p_\theta}h_\theta
=\E_p[\nabla_\theta h_\theta+h_\theta\psi].
\label{eq:measure-derivative}
\end{equation}
固定采样索引的自动微分只产生第一项。冻结测度 $q$ 下则有 $\nabla\E_qh_\theta=\E_q\nabla h_\theta$，且 $\E_q\psi$ 一般不为零。有限词表使逐项微分合法；连续空间需另行提供可交换微分与积分的支配条件。
\end{proof}

\section{散度规范化与局部二阶项}
\begin{proposition}[幂散度的生成函数与对偶]
对 $0<\kappa<1$，幂散度的规范化生成函数为
\begin{equation}
g_\kappa(x)=\frac{\kappa x+1-\kappa-x^\kappa}{\kappa(1-\kappa)}.
\end{equation}
其曲率为 $x^{\kappa-2}$；连续端点为 $g_0=x-1-\log x$、$g_1=x\log x-x+1$。交换分布时 $D_\kappa(P\Vert Q)=D_{1-\kappa}(Q\Vert P)$，且 $D_{1/2}=2H^2$。
\end{proposition}
\begin{proof}
对凸函数 $g$，定义 $D_g(P\Vert Q)=\sum_aQ_ag(P_a/Q_a)$。加入线性项 $b(x-1)$ 不改变散度。采用 $g(1)=g'(1)=0$ 的规范化后，
\begin{equation}
g_\kappa(x)=\frac{\kappa x+1-\kappa-x^\kappa}{\kappa(1-\kappa)},
\qquad g_\kappa''(x)=x^{\kappa-2}.
\end{equation}
$g_0(x)=x-1-\log x$、$g_1(x)=x\log x-x+1$ 由洛必达法则得到。对偶关系由逐点恒等式 $xg_{1-\kappa}(1/x)=g_\kappa(x)$ 给出：
\begin{equation}
D_\kappa(P\Vert Q)=D_{1-\kappa}(Q\Vert P).
\end{equation}
在 $\kappa=1/2$，$D_{1/2}(P\Vert Q)=4(1-\sum\sqrt{PQ})=2H^2(P,Q)$。
\end{proof}

\begin{proposition}[规范化散度的 Fisher 二阶项]
设 $Q>0$ 固定，$g$ 在一附近三阶连续可微，$g(1)=g'(1)=0,g''(1)=1$。若 $\sum_a\delta_a=0$ 且 $\max_a|\delta_a/Q_a|\to0$，则
\begin{equation}
D_g(Q+\delta\Vert Q)=\tfrac12\sum_a\delta_a^2/Q_a
+O\!\left(\sum_a|\delta_a|^3/Q_a^2\right).
\end{equation}
\end{proposition}
\begin{proof}
若 $g$ 在一附近三阶可微，$g''(1)=1$，并且 $P=Q+\delta$、$\sum_a\delta_a=0$、$\max_a|\delta_a/Q_a|\to0$，则
\begin{equation}
D_g(Q+\delta\Vert Q)
=\frac12\sum_a\frac{\delta_a^2}{Q_a}
+O\left(\sum_a\frac{|\delta_a|^3}{Q_a^2}\right).
\end{equation}
证明是对 $g(1+\delta_a/Q_a)$ 逐项 Taylor 展开；常数和一次项消失，余项乘以 $Q_a$ 后求和。这个共同 Fisher 二阶尺度不表示不同散度在远处相等。
\end{proof}

\section{中心化坐标与归一化恢复}
\begin{proposition}[中心化误差与概率分布的双射]
固定正分布 $q,m$。映射 $p\mapsto e=C_q\log(p/m)$ 将正单纯形双射到 $\{e:\E_qe=0\}$，逆映射为
\begin{equation}
p_a=m_a\exp[e_a-\log\E_me^e].
\end{equation}
\end{proposition}
\begin{proof}
固定 $q,m>0$。定义 $e=C_q\log(p/m)$，则 $\E_qe=0$。写 $\log(p/m)=e+c$，归一化给出 $e^c\E_me^e=1$，故
\begin{equation}
p_a=m_a e^{e_a-\Lambda(e)},\qquad
\Lambda(e)=\log\E_me^e.
\label{eq:quotient-inverse-core}
\end{equation}
反过来，任何 $\E_qe=0$ 的有限向量经该式得到正分布，并恢复原 $e$。因此归一化分布与 $q$ 中心化坐标一一对应。$\Lambda(e)$ 可以为负，因为中心化使用 $q$，配分函数使用 $m$。
\end{proof}

欧氏中心坐标 $P_\perp v$ 使用 $P_\perp=I-\mathbf1\mathbf1^\top/|\mathcal V|$，而 $C_qv$ 使用 $q$。两者相差常数，代表同一个归一化自由度；它们的欧氏范数及所在子空间不能直接混用。

\section{从对数误差到概率距离}
\begin{lemma}[加权误差控制跨度]
若 $q_i,q_j>0$，则
\begin{equation}
(v_i-v_j)^2\le(q_i^{-1}+q_j^{-1})\E_qv^2.
\end{equation}
特别地，若 $q_a\ge q_{\min}>0$，则 $(\osc v)^2\le2\E_qv^2/q_{\min}$。
\end{lemma}
\begin{proof}
将差写为 $(\sqrt{q_i}v_i,\sqrt{q_j}v_j)\cdot(q_i^{-1/2},-q_j^{-1/2})$，应用 Cauchy--Schwarz，再把两个加权平方放大为全部平方和。选择最大与最小动作即得跨度界。
\end{proof}

\begin{theorem}[Hilbert 距离到 TV 的锐界]
\label{thm:hilbert-tv}
对正概率分布，
\begin{equation}
\TV(p,q)\le\tanh\left(\frac{\osc\log(p/q)}4\right).
\end{equation}
任意给定正跨度都有二元分布取等。
\end{theorem}
\begin{proof}
令 $L=p/q$、$l=\min L$、$u=\max L$、$h=u/l=e^\delta$。$\E_qL=1$ 保证 $l\le1\le u$。凸函数 $|L-1|$ 不超过端点连线，所以
\begin{equation}
\TV(p,q)\le\frac{(u-1)(1-l)}{u-l}
=\frac{h+1-hl-l^{-1}}{h-1}
\le\frac{\sqrt h-1}{\sqrt h+1}.
\end{equation}
最后一步使用 $hl+l^{-1}\ge2\sqrt h$。取 $L$ 的两值为 $e^{\delta/2},e^{-\delta/2}$，在 $q$ 下的概率为 $1/(1+e^{\delta/2})$ 和 $e^{\delta/2}/(1+e^{\delta/2})$。它们满足均值一，且弦线及算术--几何均值不等式同时取等。
\end{proof}

\begin{theorem}[Pinsker 不等式]
对有限词表上的正概率分布，$\TV(p,q)\le\sqrt{\KL(p\Vert q)/2}$。
\end{theorem}
\begin{proof}
另一个常用界是 Pinsker 不等式 $\TV(p,q)\le\sqrt{\KL(p\Vert q)/2}$。为给出常数，取 $B=\{a:p_a\ge q_a\}$，则 $p(B)-q(B)=\TV(p,q)$。log-sum 不等式把 KL 下界为二元 KL。固定 $v=q(B)$，函数
\begin{equation}
b(u)=u\log(u/v)+(1-u)\log[(1-u)/(1-v)]
\end{equation}
满足 $b(v)=b'(v)=0$、$b''(u)=1/[u(1-u)]\ge4$。在 $v,u$ 之间积分两次，得到 $b(u)\ge2(u-v)^2$。非平凡正支持情形中 $0<v<1$；距离为零时直接成立。
\end{proof}
